\section{Introducción}

Dentro del diseño y la gestión de las redes eléctricas, la operación en paralelo de transformadores constituye un recurso fundamental para elevar la continuidad de servicio, responder de forma flexible a las variaciones de la demanda y optimizar las maniobras de mantenimiento programado. No obstante, la viabilidad técnica de esta disposición exige una estricta rigurosidad operativa, siendo imprescindible garantizar la coincidencia en las relaciones de transformación, una polaridad idéntica y una adecuada concordancia en las impedancias porcentuales de cortocircuito, evitando así la aparición de corrientes circulantes indeseadas o la sobrecarga desigual de las unidades.

El presente informe se centra en evaluar experimentalmente el acoplamiento en paralelo de unidades transformadoras monofásicas. Inicialmente, se lleva a cabo la caracterización de un transformador bajo prueba mediante la determinación de sus parámetros eléctricos fundamentales (relación de transformación, polaridad e impedancia en cortocircuitp). A partir de estos datos, se analiza el comportamiento dinámico del conjunto frente a distintas condiciones de carga, contrastando la repartición real de corriente y potencia con los modelos teóricos de circuitos equivalentes. Finalmente, el informe detalla la metodología empleada, la tabulación de los resultados de cada ensayo y una discusión analítica sobre el cumplimiento de los criterios de paralelismo.